% Figure 1: compact monochrome overview of the v2 and v3 experiment branches.
% Qualifications and measurement boundaries remain in the caption/text.
\begin{tikzpicture}[
  x=1mm,y=1mm,
  title/.style={font=\sffamily\bfseries\footnotesize,align=center},
  lab/.style={font=\sffamily\footnotesize,align=center},
  strong/.style={font=\sffamily\bfseries\footnotesize,align=center},
  box/.style={draw=black,line width=.55pt,rounded corners=.7mm,fill=white,
              minimum height=6mm,inner sep=.8mm,font=\sffamily\footnotesize,align=center},
  flow/.style={-{Latex[length=1.8mm,width=1.2mm]},line width=.6pt}
]
  \node[title] at (15,0) {(a) v2の超過};
  \node[title] at (48,0) {(b) 生存期間};
  \node[title] at (81,0) {(c) 部分保持と再計算};
  \node[title] at (116,0) {(d) RP2350};
  \node[title] at (151,0) {(e) 評価範囲};

  % (a) The v2 working set does not fit in the complete on-chip SRAM.
  \draw[line width=.65pt] (7,-31) rectangle (23,-16);
  \node[lab] at (15,-24) {SRAM\\520 KiB};
  \draw[line width=2.2pt] (15,-13) -- (15,-6);
  \draw[flow] (15,-13) -- (15,-4.5);
  \node[strong,anchor=west] at (17,-10) {6.34 MB};

  % (b) Two disjoint logical lifetimes share one physical slot.
  \draw[flow] (36,-9) -- (61,-9);
  \draw[fill=black!25,line width=.5pt] (38,-15) rectangle (48,-20);
  \node[lab] at (43,-17.5) {B};
  \draw[pattern=north east lines,pattern color=black,line width=.5pt]
    (51,-15) rectangle (61,-20);
  \node[lab,fill=white,inner sep=.3mm] at (56,-17.5) {C};
  \draw[flow] (49.5,-22) -- (49.5,-26);
  \draw[line width=.6pt] (41,-28) rectangle (58,-36);
  \draw[line width=.5pt] (41,-32) -- (58,-32);
  \node[lab] at (49.5,-30) {B};
  \node[lab] at (49.5,-34) {C};

  % (c) The v2-only retained pair returns to full precision on a tie.
  \node[box,text width=15mm] (sk) at (76,-15) {$\ell,h_{64}$};
  \node[box,text width=18mm,line width=.85pt] (ex) at (76,-30) {$N_i\ ?\ N_j$};
  \node[box,text width=9mm] (ord) at (94,-15) {順序};
  \draw[flow] (sk.east) -- node[lab,above] {$\neq$} (ord.west);
  \draw[flow] (sk.south) -- node[lab,right] {$=$} (ex.north);
  \draw[flow] (ex.east) -| (ord.south);

  % (d) v2 uses both methods; v3 reuses only the lifetime principle.
  \node[box,text width=20mm] (v2) at (116,-14) {v2提案実装\\172 kB};
  \node[box,text width=20mm] (v3) at (116,-30) {v3適応実装\\$-4.7$ kB};

  % (e) Functional and memory claims do not imply side-channel resistance.
  \node[strong] at (151,-12) {機能 $\checkmark$};
  \node[lab] at (151,-17) {メモリ $\checkmark$};
  \draw[dashed,line width=.8pt] (138,-22) -- (164,-22);
  \node[strong] at (151,-28) {定数時間 $\times$};
  \node[lab] at (151,-34) {電力・電磁波 ?};

  \draw[flow] (24.5,-22) -- (33.5,-22);
  \draw[flow] (63,-17) -- (66,-17);
  \draw[flow] (99,-15) -- (103,-15);
  % Route the lifetime-only v3 branch below the retained pair comparison box.
  \draw[flow] (63,-24) -- (64,-24) -- (64,-39) -| (v3.south);
  \draw[flow] (129,-14) -| (136,-20);
  \draw[flow] (129,-30) -| (136,-24);
\end{tikzpicture}
